% Nouveau EGA V (cote 25): archive cote 25, site page 106 (PDF page 107).
% typescript. Draft compared: none.
% AI-transcribed by Claude Opus 5.5 (claude-opus-5-5, adaptive thinking, effort high), Message Batches API, batch ega25b (msgbatch_01A7n5YuXubUPNK1zsepi8RG), 2026-11-13; not yet checked by eye.
\scanpage{25/106}
\mspage{9}

2. \uline{Diviseurs relatifs et faisceaux inversibles sur les fibrés projectifs}

\uline{et multiprojectifs}.

\hand{2.1.} Soit comme au N°1 P=P($\underline{E}$), $\underline{E}$ loc libre sur S de rang $\geqslant$ 2 partout.
Nous nous proposons de déterminer l'ensemble Div(P/S) des diviseurs rela-%
tifs \var{\inserted{$\geqslant 0$}}{$\geqslant 0$}{insertion}{ega25_25_106-v0001} sur P p. r. à S. On sait que la donnée d'un tel diviseur revient à la
donnée d'un Module inversible $\underline{L}$ sur P, muni d'une section $\phi$ de $\underline{L}$
transversalement régulière. Or d'après 1.1. (abstraction faite d'une
partition éventuelle de S, lorsque S n'est pas connexe) $\underline{L}$ est isomorphe
à un $\underline{M}\otimes\underline{O}_P(n)$, où $\underline{M}$ est un Module inversible sur S, d'ailleurs déterminé
à isomorphisme unique près en termes de $\underline{L}$. D'ailleurs, on sait que
l'on a alors des isom canoniques
\[ \hand{(*)} \qquad f_{\hand{*}}(\underline{L}) \simeq \underline{M}\otimes f_{\hand{*}}(\underline{O}_P(n)) \simeq \underline{M}\otimes\underline{Sym}^n(\underline{E}) \]
donc la donnée d'une section \var{\xout{de}}{}{typewriter strike-out}{ega25_25_106-v0002} $\phi$ de $\underline{E}$ revient à la donnée d'une section
\var{\inserted{$\phi$}}{$\phi$}{insertion}{ega25_25_106-v0003} de $\underline{M}\otimes\underline{Sym}^n$(E) . Compte tenu du fait que les fibres de P/S sont intègres,
on voit d'ailleurs que $\phi$ est transversalement régulière (i.e. régulière
sur chaque fibre) sss \hand{$\phi$}(s)$\neq$0 pour tout s$\in$S, ou ce qui revient au même,
s\hand{\unsure{ss}} l'homomorphisme
\[ \hand{\phi} : \check{\underline{M}} \longrightarrow \underline{Sym}^n(E) \]
défini par \hand{$\phi$} est ``universellement injectif'' i.e. loc un isomorphisme sur
un facteur direct, ou ce qui revient au même, si son transposé
\[ \hand{\check{\phi}} : \underline{Sym}^n(\underline{E})^{\hand{\vee}} \longrightarrow \underline{M} \]
est surjectif.

Nous dirons que le diviseur relatif D sur P est \uline{de degré} n
si $\underline{O}_X(\underline{D})= \underline{L}$ est de degré n au sens du n° précédent. \var{\deleted{Ainsi, grâce à 1.1.}}{}{deletion}{ega25_25_106-v0004}
\var{\deleted{pour D donné, il existe une unique \var{\inserted{décomposition}}{décomposition}{insertion}{ega25_25_106-v0006} \var{\deleted{partition}}{}{deletion}{ega25_25_106-v0007} de S en somme d'ouverts $S_n$}}{}{deletion}{ega25_25_106-v0005}
\var{\deleted{(n N) tel}}{}{deletion}{ega25_25_106-v0008}

\mspage{106}
